\section{Design}\label{sec:design}

The library has three layers and a verb. \emph{Algebra} defines cells, chains and coefficient fields. \emph{Streams} produce
filtered complexes as sequences of cells in filtration order. \emph{Engines} consume a stream and return a barcode with
representatives. \code{Persistence} chooses a stream, an engine and defaults from one call.

\subsection{Cells, chains and fields}\label{sec:algebra}

\paragraph{Fields are typeclasses.}
A coefficient type is a type with a \code{Field} instance. \code{FiniteField(p)} gives the field with $p$ elements; its
element type \code{Fp} is an \emph{opaque type per prime}, so elements of $\mathbb{F}_5$ and $\mathbb{F}_7$ have
incompatible types and mixing them is a compile-time error. Exact arithmetic over $\Q$ is available for small examples;
double-precision coefficients take a tolerance. Engines never inspect the runtime class of a coefficient: on Scala.js every
boxed number is a \code{Double}, so a test such as \code{case d: Double} would also match an \code{Fp}. They ask the field.

\paragraph{Cells carry their boundary.}
A \code{Cell} has a dimension and a \code{boundary}, a list of faces with coefficients. An \code{OrderedCell} adds a total
order. The library's instances are simplices (on any vertex type), cubes (in doubled coordinates, so that a cube of any
dimension is a tuple of integers), the generators of a finite simplicial set, and cells of a quotient. A new kind of cell is a
new instance, and every generic engine runs on it. There is deliberately no dual notion of a cocell: a coboundary is
extrinsic to a cell, since it depends on which higher cells the ambient complex contains, so engines that need
coboundaries compute them against the complex they are given.

\paragraph{Chains are formal sums with pluggable storage.}
A \code{Chain[CellT, CoefficientT]} is a formal sum of cells, kept in the order of the cells; its leading term is the
reduction pivot. \code{Chain} is abstract and open. Equality and hash are those of the formal sum, whatever the storage, so a
chain stored as a heap, packed into primitive arrays, or expanded lazily from a reduction log compares equal to the same sum
stored any other way. The listing shows the style, which uses Scala~3's context-bound syntax for typeclasses.

\begin{lstlisting}[language=Scala]
trait Cell extends HasDimension:
  type Self
  extension (self: Self)
    def boundary[CoefficientT: Field]: Seq[(Self, CoefficientT)]

trait OrderedCell extends Cell:
  type Self: Ordering as ordering
\end{lstlisting}

\paragraph{Default instances live with the data, never with the typeclass.}
The core package has no top-level \code{given}. The default instances for \code{Simplex}, \code{Cube}, \code{Chain} and the
other data types sit in their companions, which puts them in the types' implicit scope: found without an import, and
consulted only when nothing lexically visible matches, so a user's own instance never ties with them. An instance in the
companion of a \emph{typeclass} (\code{Field}, \code{OrderedCell}) would be searched while the instance type was still
unknown and would silently decide type inference: a default \code{Double is Field} would turn a forgotten import of a
finite field into real coefficients and a wrong torsion answer. A spec checks the rule from a user's side, with code
that sees only the public imports. The price is an import line, \code{import scala.language.experimental.modularity},
in every user file, because the typeclass syntax is experimental in Scala~3.9.
\VERIFY{whether the import is still required in the release the paper describes}

\subsection{Filtered complexes are streams in filtration order}\label{sec:streams}

A \emph{stream} yields cells in filtration order, together with each cell's value. \tdaj{} requires of every stream the
same \emph{ordering contract}. The order is one total order, built by one combinator: filtration value, then dimension, then
a tie-break (for Vietoris--Rips the colexicographic order of the combinatorial index, as in Ripser). Within a dimension the
cells are delivered in that same order, never in an independently built one that agrees except on ties. Values are monotone
exactly: a face is never larger than its coface, down to the last bit of a floating-point value. Engines pair under the
stream's order, so a stream whose ties disagree with the engine that reads it produces a barcode that is wrong in a way the
Betti numbers do not show. This has been the largest single source of bugs in the library, and the contract is tested cell
for cell and value for value against an independent construction (\cref{sec:oracles}).

Streams are level-wise: a Vietoris--Rips stream generates the simplices of one dimension at a time, on flat arrays, and
an engine reading it up to degree $k$ never forces dimension $k+2$. A stream says which homological degrees it covers
(\code{homologyDegreeLimit}), and every engine that reads a stream stops there, so the artifacts at the cut of a truncated
complex are never reported as classes. The dispatchers (\code{VietorisRips}, \code{Cech}, \code{Witness},
\code{Dowker}, \code{DtmRips}, \code{SparseRips}, \code{Truncated}, \code{AlphaShapes}, \code{CubicalImage}) take \code{maxDimension} as the top
\emph{homological} degree, in contrast to JavaPlex and GUDHI where it is the top simplex dimension; the construction classes
behind them are public machinery, and \cref{app:catalogue} lists them.

\subsection{Engines and the verb}\label{sec:verb}

Engines are independent implementations that share \code{Chain} primitives (\cref{tab:engines}); a fix to one says nothing about
the others, and each fast engine has a slower reference in the test tree that it must equal term for term.

\begin{table}[t]
\centering\small
\begin{tabular}{@{}p{3.9cm} p{3.4cm} p{3.2cm} p{3.8cm}@{}}
\toprule
Engine & Input & Native kind & Role \\
\midrule
Naive & any ordered cell & cycles & reference; cursor \\
Chunks & any stream & cycles & clear-and-compress~\cite{bkr2014chunks} \\
Cohomology & any stream & cocycles & generic, dense numbering \\
Ripser (packed) & Rips of a metric space & cocycles & the default for Rips \\
Parallel Ripser & Rips of a metric space & cocycles (lazy) & edge blocks on threads \\
Packed cubical & cubical grid & cocycles & image cocycles \\
Fast cubical & cubical grid & cycles & dual union-find; default for images \\
Fast alpha & Delaunay alpha complex & cycles & dual union-find \\
\bottomrule
\end{tabular}
\caption{The engines behind \code{Persistence}. Each returns the other kind of representative by the involution
(\cref{sec:representatives}), except the two fast engines (cycles only). A tenth class, the Ripser engine on explicit
\code{Simplex[Int]} cells, is kept as a test oracle and a block-parallel experiment on explicit matrices is kept for its
helpers; neither is reachable from the verb.}
\label{tab:engines}
\end{table}

The verb is one call:
\begin{lstlisting}[language=Scala]
import scala.language.experimental.modularity
import org.appliedtopology.tda4j.*

val points = Array.tabulate(30)(i => Array(math.cos(i * 0.21), math.sin(i * 0.21)))
val diagram = Persistence(points)   // degrees 0 to 2, coefficients in the field with 17 elements
diagram.dim(1).longest              // the most persistent loop, with a representative cycle
\end{lstlisting}
Its full form is \code{Persistence(input, maxDimension = 2, maxFiltrationValue, complex = VietorisRips, characteristic = 17,
engine = Auto)}. The input is anything that converts to a stream: points, a metric space, an image, or a stream. Public numeric
options accept \code{2.0}, \code{2}, \code{Some(2.0)} or \code{None}, and point inputs accept the array and sequence types
users have; every such parameter is an \code{into} type so that a single \code{apply} with defaults covers the call.
\code{Engine.Auto} chooses the packed Ripser engine, or the parallel one when several threads are available and it applies,
for Vietoris--Rips complexes on points or a metric space; the fast cubical engine for images of dimension two and up when
cycles are asked for; and the generic cohomology engine otherwise.
The result is an immutable \code{PersistenceDiagram}: bars with representatives, the coefficient field as a type member,
and queries (\code{dim}, \code{at(f)}, \code{longest}, \code{longerThan}, \code{bettiNumbers}).

\paragraph{The cursor.}
Where the run is long, an engine's cursor lets the caller advance it for a time slice, read progress
(\code{processedCells}, \code{totalCells}), and ask for \code{diagramAt(f)}, which is exact at any filtration value $f$
wherever the cursor stands. Only the homology engines are prefix-ordered: the naive engine has the cursor, and the chunks
engine answers the same queries without one. The cohomology engines, which are the fast ones, process cells in a different
order and cannot offer it.

\subsection{Defaults that encode lessons}\label{sec:defaults}

\begin{itemize}[leftmargin=1.5em, itemsep=1pt]
\item \textbf{Coefficients are $\Fsev$.} $\Ftwo$ hides signs and odd torsion: an engine with a sign error passes every
  $\Ftwo$ test. The default is also the default of the MATLAB and command-line front ends, so a user who wants to compare
  with Ripser or GUDHI passes \code{characteristic = 2} on purpose.
\item \textbf{Zero-length bars are dropped.} Every engine, the verb and the front ends take \code{includeZeroLength},
  default \code{false}. A test must not use ``number of bars equals number of cells'' as an oracle unless it opts in.
\item \textbf{Degrees $0$ to $2$ are computed.} Higher degrees cost far more. The verb caps the reduction at the requested
  degree, and a truncated stream carries its own limit (\cref{sec:streams}).
\item \textbf{The threshold is the enclosing radius.} For a point cloud the enclosing radius, the smallest over the points
  of the largest distance to another point, is where the complex becomes a cone and the barcode stops changing. It
  is Ripser's default as well, and credited there to Henselman-Petrusek~\cite{bauer2021ripser}.
\end{itemize}

\subsection{Platforms and interfaces}\label{sec:platforms}

\paragraph{One source tree, two platforms.}
The core compiles from one tree for the JVM and, through Scala.js, for JavaScript; platform-specific code lives behind a
handful of seams (parallelism, file access, the live plot viewer, the image decoder, the filtration cache), each defined
once per platform under the same name. Shared code uses no Java library that Scala.js lacks, which is why the library
carries its own Jacobi singular-value decomposition, linear solvers, vantage-point tree and enclosing-ball solver.
Parallel steps run on a fork-join pool on the JVM and sequentially in JavaScript. Specs run on both platforms; a property
that searches for rare inputs runs all its cases on the JVM and a tenth of them on Scala.js, since only what differs between
the platforms needs the second run.

\paragraph{Front ends.}
The JVM library is called directly from Scala and from Java. A facade of plain arrays and option strings
(\code{matlab.TDA4j}) is the entry point for MATLAB, Java and other JVM languages; the command-line program is a thin
translator over it, and the Python package (a wheel over the facade via JPype) and the npm package (ES modules and type
declarations over the Scala.js build) mirror it. Adding an option means adding it to the facade, the command line and the
Python package together; a test fails when the Python package falls behind.
\VERIFY{the MATLAB facade has not been run in MATLAB as of PLAN-paper.md 2026-10-04; its bytecode needs a MATLAB release with Java~17 or later (R2024a on MathWorks' compatibility page, read through a summarizer only)}
\VERIFY{Python package: published on PyPI or not at submission (not yet at 2026-11-01)}

\paragraph{Documentation is tested.}
The code fences of the tutorials and the user guide are compiled by the documentation build, and the numbers the tutorials
quote are asserted by specs, so the documentation cannot drift from the library without a build failing; the same example
appears in each language as tabs.
\VERIFY{the exact coverage: fences marked as not compiled exist, and the Python and MATLAB tabs are checked by other means} \Cref{app:first-lines} compares the first lines a newcomer writes for ``barcode with representatives from a point cloud'' in
\tdaj{} and in the tools of \cref{tab:related}.

\subsection{Openness}\label{sec:openness}

The library is a platform for experiments, so public and extensible is the default. Access modifiers are not documentation:
``this is not the entry point'' is said in the scaladoc, not by hiding the class. A restriction carries a comment naming one
of five reasons: invariant-bearing mutable state; a library-wide guarantee (nothing public may drop representatives); the
flat namespace (a public top-level name joins every user's wildcard import and shadows their own definitions); a facade whose
contract is strings and arrays; or a name that is public another way. Hooks are kept even where no current entry point uses
them: a \code{keepCriterion} on the Vietoris--Rips generator lets an experiment prune simplices, with the caveat, documented
rather than checked at run time, that the kept set must be closed under taking faces.
